\section{Discussion}

We developed a fully automated pipeline for scratch-wound quantification that achieved 92.1\% end-to-end expert-validated accuracy across 126 publication-condition phase-contrast images (116/126, with the automated QC decision treated as part of the pipeline output) and 97.2\% boundary-detection accuracy on the 106 wound-bearing images (103/106) --- all without annotated training data or deep-learning models. Like existing computational approaches, the method applies learned parameters (local variance window size, edge thresholds, smoothing window lengths) that were developed on this dataset; full parameter documentation is provided to support transparency and reproducibility.

As expected for frame-wise segmentation, the unconstrained baseline produced substantial non-monotonic area fluctuations (violation rate 44.7\%, MASD 12003~$\mu$m$^2$). The pipeline integrates temporal constraints into the segmentation stage, coupling texture-based single-frame detection with a per-column Kalman filter and hard monotonic closure constraint. This hierarchical design allows the Kalman filter to down-weight artifact-contaminated observations while the hard constraint provides a correctness guarantee. The resulting trajectories showed monotonically decreasing wound area and substantially lower frame-to-frame variance (ablation arm D: MASD 1579~$\mu$m$^2$), representing a 96.8\% improvement over unconstrained detection. The pipeline is released as open-source software at \url{https://github.com/GiorgioRicciardiello/WoundDetectionBounds} \citep{Ricciardiello2025code}, requiring only NumPy, OpenCV, and SciPy on standard laboratory hardware.

The single-frame detector relies on a texture-based observation: cells scatter light and produce high local variance in phase-contrast images, while the cell-free wound channel does not \citep{Yin2012}. This variance contrast varies across imaging conditions; training-free methods may be sensitive to changes in magnification, illumination, or microscope settings. The detection accuracy achieved without annotated training data was 92.1\% end-to-end and 97.2\% on wound-bearing images, falling within the range reported for supervised methods on comparable assays, although direct comparison on identical images was not performed. The two-stage boundary regularization — RANSAC-inspired outlier rejection followed by Savitzky--Golay smoothing — reduced edge artifacts; the relative contributions of outlier rejection versus smoothing to the final boundary accuracy are not independently quantified \citep{Fischler1981,Savitzky1964}.

The monotonic closure constraint encodes an assumption about the assay: wound area should not increase over time in a typical scratch-wound experiment. This assumption is valid when closure dynamics dominate over transient detector noise; it does not apply when genuine biological re-opening occurs (e.g., cytotoxic concentrations, cell retraction during early phases). Enforcing this per column — stronger than a global area constraint — reduced the frame-to-frame violation rate from 44.7\% (unconstrained) to 0\% (constrained). The Kalman-filtered formulation combines this constraint with observation weighting: rather than clamping edges deterministically, the filter down-weights observations with high residual variance arising from debris or illumination artifacts, while the hard clamp acts as a backstop to enforce the monotonicity assumption on every frame. The extracted trajectories exhibited monotonically decreasing wound area; a linear mixed-effects model applied to these trajectories yielded a significant closure-over-time main effect ($\beta = 0.959$~$\mu$m/h, $p = 9.3 \times 10^{-6}$; 95\%~CI~[0.534, 1.383]; Equation~\eqref{eq:mixed-model}). Cross-sectional pairwise comparisons between conditions yielded corrected $p$-values of 1.000; given the modest closure range (${\sim}17\%$ in 20~h), small per-group sample sizes ($n = 5$--28), and this test-bed design, larger effect sizes would be required for pairwise significance.

The automated quality-control module achieved sensitivity of 1.00 (all 103 wound-bearing images with correct boundaries were accepted; zero false negatives) and specificity of 0.565 (Cohen's $\kappa = 0.68$) across the 126 publication-condition images. The lower specificity reflects a deliberate design choice: avoiding false negatives (permanent data loss) was prioritized over minimizing false positives on no-wound fields, since the latter are excluded from dynamics analysis anyway. Of the 23 no-wound fields, 7 produced spurious masks; expert review identified all 7.

However, this design point is dataset-specific. \textbf{Extension to new imaging conditions, microscope settings, cell types, or magnifications will likely require re-tuning of the QC thresholds.} The validation criteria (aspect ratio, edge homogeneity, width coverage) were optimized empirically on this 126-image cohort; their performance on new datasets remains to be established. We provide source code and threshold documentation to support this process; users are advised to validate on a small local cohort (e.g., 20--30 manually annotated images) before applying the pipeline to large datasets in a new context.

Of the 10 expert--detected errors, 7 were no-wound fields where the detector produced spurious masks, and 3 were wound-bearing images with minor boundary-segment localisation errors under strict annotation criteria (eight of the 10 concentrated in the Media condition; the remaining two in DMSO). The 3 boundary-shift errors occurred in wound-bearing images where the wound channel was identified but boundary segments in a subset of columns deviated from the manual trace. The review was conducted under strict criteria --- any visible deviation in boundary alignment, even if spatially localised, was classified as incorrect. This strict criterion was applied to establish a conservative accuracy bound; looser annotation criteria would yield higher accuracy estimates. The 7 spurious masks on no-wound fields were produced by the detector but not rejected by QC; these 7 trajectories were excluded from the dynamics analysis, which used only images where both detector and QC agreed on wound presence.

\textbf{Validation was performed by a single trained observer, blinded to treatment condition but not cell-line identity.} While the strict annotation criterion renders the reported accuracy a conservative lower bound, single-observer validation is an important limitation. Inter-rater reliability was not quantified, and we cannot exclude systematic biases in the reviewer's threshold or condition-dependent drift in decision-making. For high-stakes applications, independent validation on domain-specific datasets before deployment is recommended. Future work should include multi-observer validation (inter-rater ICC or Cohen's $\kappa$) and testing on imaging systems or cell types not present in this cohort.

The ablation study (Section~\ref{sec:ablation-design}) provides an empirical decomposition of the two temporal constraint components (Table~\ref{tab:ablation}).
The hard clamp alone (arm~B) drove the violation rate from 44.7\% to zero and substantially reduced frame-to-frame variance ($\sigma^2_\text{FTF} = 11.7 \times 10^6$ vs.\ $206 \times 10^6$~$\mu$m$^4$ at baseline), but produced the roughest constrained trajectories (MASD $= 1796$~$\mu$m$^2$, 95\%~CI 1604--2036).
The Kalman filter is the primary driver of smoothness: Kalman-only (arm~C) achieved MASD of $1230$~$\mu$m$^2$ (95\%~CI 1022--1544), a reduction that the hard clamp alone cannot replicate.
Arm~C produced zero observed violations in this dataset; however, the Kalman filter provides no formal monotonicity guarantee --- on frames where the posterior estimate marginally exceeds the prior wound area, the area may increase slightly --- so Kalman-only cannot serve as a correctness guarantee in the general case.
The Kalman+hard production configuration (arm~D) combines Kalman smoothing with hard clamping.
Arm~D achieved MASD of $1579$~$\mu$m$^2$ (95\%~CI 1336--1934), compared to arm~C (Kalman-only) MASD of $1230$~$\mu$m$^2$ (95\%~CI 1022--1544). Computed as $(12003 - 1579)/(12003 - 1230) = 96.8\%$, arm~D retained $96.8\%$ of the smoothness gain of arm~C relative to unconstrained (arm~A), where $12003$~$\mu$m$^2$ is the unconstrained baseline MASD. Arm~D produced zero observed frame-level monotonicity violations; arm~C also produced zero observed violations, although the Kalman filter provides no formal monotonicity guarantee.
The increased roughness of arm~D relative to arm~C (MASD $1579$ vs.\ $1230$~$\mu$m$^2$) corresponds to frames where hard clamping was applied. At the per-frame level, hard clamping occurred on at least one column in essentially every frame (per-frame, per-trajectory constraint rate $= 100\%$; Supplementary Table~\ref{tab:supp-qc-validation}). The aggregate wound-area trajectory in arm~D closely tracked arm~B (hard-clamp only), with closure fractions differing by only 0.015 (arm~B: 0.191 vs.\ arm~D: 0.176).

Closure fractions at the final timepoint differed across the four arms. Unconstrained (arm~A) yielded 0.021 (95\%~CI $-0.009$--0.055), hard-clamp only (arm~B) 0.191 (95\%~CI 0.173--0.215), Kalman-only (arm~C) 0.086 (95\%~CI 0.069--0.110), and Kalman+hard (arm~D) 0.176 (95\%~CI 0.159--0.198). Arm~C (Kalman-only) achieved the lowest trajectory roughness (MASD $1230$~$\mu$m$^2$) but reported lower closure fraction than arm~D (0.086 vs.\ 0.176). Arm~D reported closure fraction intermediate between arm~C and arm~B, differing from arm~B by 0.015.

The approach has limitations that should be acknowledged. Significant departure from normality was observed in the Line~2 DMSO and Line~2 Media groups (Shapiro--Wilk $p < 0.001$; Supplementary Table~\ref{tab:supp-assumption-checks}); while mixed-effects fixed-effect estimates are generally robust to moderate non-normality, the cross-sectional Welch's $t$-tests for Line~2 ALK5i ($n = 5$) should be interpreted with caution given the small sample and non-normal reference distribution. The monotonic constraint carries an irreversibility cost at the per-column level: a sufficiently large spurious inward detection, from a bubble or severe focus shift, will propagate to later frames even under Kalman filtering. The edge-detection threshold assumes unimodal cell variance and may require adaptation for heterogeneous populations or images with illumination gradients. Edges are resolved at integer-pixel resolution, introducing quantization error as wounds approach closure; measurements are reported in calibrated physical units ($1.24\,\mu$m/pixel), but the underlying spatial resolution remains limited by the pixel grid. A larger expert-annotated dataset would further tighten estimates of detection accuracy across imaging conditions and cell types. Detection accuracy and quality-control performance were established against a single trained observer: although the reviewer was blinded to treatment condition, cell-line identity was not masked (the lines are morphologically distinguishable), and neither a second independent observer nor a repeated within-observer review was performed, so inter-rater reliability was not quantified. The deliberately strict annotation criterion --- any visible boundary deviation classified as incorrect --- renders the reported accuracy a conservative lower bound; nonetheless, the point estimate could shift with additional annotators, and any condition- or line-dependent drift in a single reviewer's threshold would not be detectable from these data.

Future development will address sub-pixel edge localization via parabolic fitting, adaptive per-column thresholding, and bootstrap-based uncertainty bounds for area measurements. The strictly monotonic, smooth trajectories produced by this pipeline are also well-conditioned inputs for fitting mechanistic models of collective migration — in particular Fisher--Kolmogorov reaction-diffusion frameworks \citep{Jin2015,Maini2020,Simpson2013} — enabling direct estimation of cell motility and proliferation rates from the same trajectory data.

On this test bed, the pipeline produced wound-area trajectories with monotonically decreasing area across 126 publication-condition phase-contrast images spanning two cell lines, with 92.1\% end-to-end expert-validated accuracy (116/126) and 97.2\% boundary-detection accuracy on the 106 wound-bearing images (103/106 under strict annotation criteria). The 10 detection errors comprised 7 spurious masks on no-wound fields and 3 boundary-segment shifts in wound-bearing images; all were identified at expert review. The QC module accepted 7 of the 10 erroneous detections (the spurious no-wound masks). The 106 QC-validated trajectories exhibited monotonically decreasing area and supported statistical modeling with a significant time effect ($\beta = 0.959$~$\mu$m/h, $p = 9.3 \times 10^{-6}$). The pipeline requires no annotated training data and runs in linear time per frame on standard laboratory hardware. The full implementation is available at \url{https://github.com/GiorgioRicciardiello/WoundDetectionBounds} \citep{Ricciardiello2025code}.
